% TOP_PROBLEM: {"id": 3000022, "title": "Covering a crossing supermodular function with pairwise non-parallel arcs", "category_id": 2, "category": {"id": 2, "name": "combinatorics", "display_name": "Combinatorics", "description": "Counting problems, graph theory, discrete structures.", "slug": "combinatorics", "order_index": 2, "created_at": "2026-07-31T15:26:25.671Z"}, "status": "open", "problem_number": "AMR-029-0022", "set_id": 13}
% FIRST_POSTED: 2026-10-01
% ENTRY_KIND: statement_only
% UnsolvedMath ID 3000022; source catalogue code AMR-029-0022
% !TeX program = lualatex
% Definition and sources only; this file is not an LLM solution attempt.
\documentclass[11pt,a4paper]{article}
\usepackage[margin=25mm]{geometry}
\usepackage{fontspec}
\setmainfont{lmroman10-regular.otf}[BoldFont=lmroman10-bold.otf,
 ItalicFont=lmroman10-italic.otf,BoldItalicFont=lmroman10-bolditalic.otf]
\usepackage{amsmath,amssymb,mathtools}
\usepackage{unicode-math}
\setmathfont{latinmodern-math.otf}
\AtBeginDocument{\let\setminus\smallsetminus}
\usepackage{xurl,hyperref}
\hypersetup{colorlinks=true,urlcolor=blue}
\setcounter{secnumdepth}{0}
\setlength{\parindent}{0pt}
\setlength{\parskip}{5pt}
\setlength{\emergencystretch}{3em}
\begin{document}
\section*{Covering a crossing supermodular function with pairwise non-parallel arcs}\label{3000022}
{\small UnsolvedMath ID 3000022 · AMR-029-0022 · Combinatorics · AMR Open Problem Lists}
\subsection{Definitions and mathematical statement}
Let $V$ be finite and $p:2^V\to\mathbb Z$ satisfy
$p(\varnothing)=p(V)=0$. Call $p$ crossing supermodular if
\[
                 p(X)+p(Y)\le p(X\cap Y)+p(X\cup Y)
\]
whenever $X\cap Y$, $X\setminus Y$, $Y\setminus X$, and
$V\setminus(X\cup Y)$ are all nonempty.
For a loopless digraph $D=(V,A)$, let $\rho_D(X)$ count arcs whose
tail lies outside $X$ and whose head lies inside $X$.

Determine the minimum size of an arc set satisfying
\[
 \rho_D(X)\ge p(X)\quad(X\subseteq V),\qquad
 A\subseteq\{(u,v):u,v\in V,\ u\ne v\}.
\]
Thus each ordered pair may be used at most once, while the two opposite
arcs $(u,v)$ and $(v,u)$ may both occur. Seek an effective feasibility
criterion, a characterization of the optimal value, and an attaining
construction; assign value $\infty$ when no feasible digraph exists.

The prohibition on parallel arcs is central. Unrestricted multiarc
coverings obey different conditions and may meet demands that no
simple digraph can meet. For example, a necessary condition here is
$p(X)\le |X|\,|V\setminus X|$ for every $X$, but the problem asks for
the complete interaction of all cut requirements, not just their
individual upper bounds.
\subsection{Short English statement}
How many distinct directed arcs are needed to cover all demands
of a crossing-supermodular function when repeated parallel arcs are forbidden?
\subsection{Sources}
\begin{itemize}
\item[S1] ULAM.AI and the UnsolvedMath Contributors, supplied problems.json, record 3000022 (AMR-029-0022), AMR Open Problem Lists. Catalogue identity and source transcription; CC BY 4.0.
\url{https://huggingface.co/datasets/ulamai/UnsolvedMath}.
\item[S2] Egres Open - List of open problems, Open-problem page: Covering a crossing supermodular function with pairwise non-parallel arcs. Original source indexed by AMR-029-0022.
\url{https://oldlemon.cs.elte.hu/egres/open/List_of_open_problems}.
\item[S3] EGRES Open, Covering a crossing supermodular function with pairwise non-parallel arcs. Individual source page and explanatory remarks.
\url{https://lemon.cs.elte.hu/egres/open/Covering_a_crossing_supermodular_function_with_pairwise_non-parallel_arcs}.
\end{itemize}
\end{document}
